\section{Threat model and trust roots}

\subsection{Objects and principals}

Let a campaign contain requests $q\in Q$, execution attempts $x\in X$, receipts
$r\in R$, evidence objects $e\in E$ and adjudications $a\in A$. An executor
principal is an authenticated subject allowed to perform an attempt. A provider is
the service or runtime that performs a model-mediated operation, and a checkpoint
is the immutable model identity when such an operation is used. A verifier checks
receipt structure and cryptographic or content bindings. An adjudicator applies a
frozen scientific protocol to admissible evidence. These roles may be implemented
by one organization, but their logical separation must remain explicit.

The protected objects are not only final files. They include the protocol, prompts,
normalization rules, scorer, code revision, tool versions, input and evidence
manifests, environment image, random seeds, session state and every request/result
binding. Supply-chain systems similarly distinguish an artifact from a signed
statement about the steps and principals that produced it
\citep{torresarias2019intoto,slsa2026provenance}. Our threat model applies this
distinction to research execution while declining to infer scientific validity from
software integrity.

\subsection{Adversary capabilities}

The adversary may:

\begin{enumerate}[label=\textbf{T\arabic*.}]
  \item replace an executor or model while copying a declared identity;
  \item replay a valid result under a different request, epoch or evidence set;
  \item let one nominal lane read the other lane's prompt, state or output;
  \item submit the same model or service twice under different lane names;
  \item write a success marker before all output bytes are durable;
  \item provide a non-empty but unauthenticated certificate identifier;
  \item convert a host, network, quota or tool failure into a scientific datum;
  \item omit inconvenient failures or downgrade them to missing rows;
  \item change normalization, scoring or stopping after observing an outcome;
  \item exploit a shared helper, gold function or evidence cache to manufacture agreement.
\end{enumerate}

The adversary need not break a hash. Correct hashes can bind the wrong semantics,
and signed statements can be authentic statements from an untrusted or
inappropriately scoped issuer. Identity-oriented signing and transparency logs
illustrate the additional trust structure required beyond a digest
\citep{sigstore2026signing}.

\subsection{Out-of-scope threats}

The contract does not decide whether a scientific source is correct, whether a
model response is well reasoned, whether an evaluator's ontology is complete, or
whether a statistical estimand answers the substantive question. It also does not
protect a fully compromised trust root or guarantee confidentiality. These are
separate evidence-quality, scientific-design and security questions. A deployment
may attach additional controls, but must not advertise them as properties of the
base receipt contract.

